\section{Введение}

ИИ-агент работает циклом. Модель читает накопленный контекст, вызывает инструмент, получает его вывод и снова
читает контекст, к которому добавилась новая порция: файл, страница поиска, лог тестов, ответ базы данных. Так
работают агенты программистов, исследователей и аналитиков, везде, где агент сам ищет, читает и проверяет. На машине с
унифицированной памятью такой цикл упирается в чтение. Apple M-серия, AMD Strix Halo и NVIDIA DGX Spark вмещают модель
класса 27B и контекст в сотни тысяч токенов и генерируют с приемлемой скоростью, но вычислительных блоков у них в разы
меньше, чем у дискретных GPU. Мы измерили это на M3 Max с моделью Qwen3.8-27B (квантизация Q4\_K\_M, llama.cpp на
Metal): на глубине от 131 до 181 тысячи токенов чтение идёт со скоростью от 41 до 28 токенов в секунду. Порция в
10 тысяч токенов ждёт от четырёх до шести минут, и всё это время агент стоит.

Облачный API снимает это ожидание: одна B300 при загрузке вычислителей 30\% читает такую порцию примерно за
секунду (наша оценка по пиковой производительности FP8). Поставщик
хранит кэш контекста между вызовами, но берёт плату за его чтение на каждом вызове. Это налог на поддержание кэша, и в агентной работе на глубоком контексте у медианного поставщика на него приходится большая часть расходов. На нашем профиле средняя глубина контекста 164 тысячи токенов, а новых токенов в вызове около тысячи. Даже со
скидкой на чтение кэша вызов перечитывает примерно в 150 раз больше токенов, чем добавляет: у медианного поставщика Qwen3.8-27B (\$0.05 за миллион токенов чтения кэша,
\$0.18 за миллион новых и \$2.30 за миллион выходных) чтение кэша даёт около 85\% цены вызова, у самого дешёвого около 48\%. Кроме того, облачный инференс требует полного доверия: клиент не может проверить, что сервер
сделал с его контекстом.

Мы задаёмся вопросом, можно ли оставить сессию и генерацию на машине пользователя, а чтение контекста отдать серверу,
и сколько это стоит при настоящей агентной нагрузке. Рассматриваем две модели со слоями линейного внимания:
плотную Qwen3.8-27B \cite{qwen38} и MoE Qwen3.8-Flash-Next \cite{qwen38flashnext}. Третью модель, MoE GLM-5.3-Flash \cite{glm53flash}, мы берём для тех, кто может позволить себе
машину мощнее, например Mac Studio с 256~ГБ, и закрывать через гибрид более широкий круг задач. Мы разбираем архитектуру и экономику схемы;
качество моделей не оцениваем: проверок качества мы не делали и проверяем только работоспособность самой схемы.

Схема устроена так. Сервер держит копию каждой сессии и считает prefill только для новых токенов. Клиенту он
отправляет память модели для новых позиций: KV-кэш, то есть ключи и значения слоёв внимания, и рекуррентное
состояние линейных слоёв. Клиент продолжает генерацию так, как если бы прочитал порцию сам. Большая часть ожидания prefill уходит из агентского цикла: крупные порции сервер читает за секунды, а малые
клиент дочитывает сам. Генерацию ограничивает полоса памяти, и
параллельные агенты используют её лучше: на M3 Max четыре потока на глубине 32 тысячи дали в 1.56 раза больше токенов в секунду, чем один,
по времени запросов.

Работа описывает, что передаётся между клиентом и сервером и почему (раздел~\ref{sec:architecture}), какой для этого
нужен канал (раздел~\ref{sec:channel}), каким клиентам схема выгоднее (раздел~\ref{sec:clients}), как устроена
серверная сторона (раздел~\ref{sec:server}), во что обходится схема по сравнению с API и своим серверным кластером для
команд от 5 до 5000 человек (раздел~\ref{sec:economics}), что даёт контекст в миллион токенов
(раздел~\ref{sec:million}) и что меняется в доверии к поставщику (раздел~\ref{sec:security}).

\textbf{Место работы.} Разнесение prefill и генерации давно стало стандартом внутри дата-центров
\cite{zhong2024distserve,patel2023splitwise,qin2024mooncake}, и в 2026 году оно начало выходить за их пределы. Длинный
prefill переносят на отдельные кластеры и передают кэш в другой дата-центр по обычному Ethernet \cite{qin2026prfaas};
облачный KV-кэш переиспользуемого контекста загружают на устройство \cite{liu2026sparkv}; в многоходовых диалогах
выбирают, где читать добавочный prefill, на узле prefill или на узле генерации \cite{li2026ppd}; на нескольких машинах prefill распределяют по слоям и затем ведут генерацию на одной из них \cite{ds4pr401}. Перенос кэша между разными моделями изучает NVIDIA \cite{heo2026crossmodel}, обмен латентными
представлениями между моделями объявила компания Mostik \cite{mostik2026}. Мы переносим разнесение prefill и генерации одной модели
на границу между поставщиком и персональной машиной пользователя и рассматриваем его как услугу для длительных
агентских сессий. Наш вклад:
\begin{itemize}
\item протокол двух копий растущей сессии у поставщика и у клиента: вверх идут идентификаторы токенов, вниз KV-кэш
новых позиций и рекуррентное состояние, журнал и контрольные точки держат копии согласованными, а решение отдать порцию
серверу клиент принимает на каждом вызове;
\item профиль агентской нагрузки из 1.5 миллиона вызовов и воспроизводимый движок расчётов, который показывает, при
каком канале, глубине контекста, доле открытой модели и классе клиента схема выгодна;
\item экономика для отдельного специалиста, команды и закрытого контура: аренда и покупка, весовые категории машин
на 64, 128 и 256~ГБ, контекст до миллиона токенов;
\item модель доверия: что клиент может проверить у поставщика, которому не обязан доверять, и где проверка
заканчивается;
\item прототип на llama.cpp, на котором базовый обмен проверен на процессоре с малой моделью
(приложение~\ref{app:prototype}).
\end{itemize}
